\section{Reproducibility and availability}
\label{app:reproducibility}
To regenerate the tables, install Node 20 or newer and run \texttt{node scripts/collect.mjs} from a checkout of this repository. The script's only network use is \texttt{npx} resolving the packages named in \texttt{paper/data/suite.txt}. To rebuild the PDF, install Tectonic 0.17.0 and run \texttt{make pdf}; \texttt{make arxiv} builds the submission packet and recompiles it flat. The bootstrap seed and the resample count are the engine defaults unless a caller overrides them. Public task generators use seed \suitePublicSeed{} where the provenance command prints \texttt{-{}-seed}. The defect counts in Section~\ref{sec:review} regenerate from \texttt{paper/data/defects.csv}, the builder and reviewer counts from \texttt{paper/data/builders.csv} and \texttt{paper/data/reviewers.csv} (rebuilt by \texttt{scripts/builders.mjs}), the council counts from \texttt{paper/data/council.csv} (rebuilt by \texttt{scripts/council.mjs}), the fix-round counts from \texttt{paper/data/review-rounds.csv}, and the preregistered trivial-policy scores from \texttt{paper/data/trivial.csv}; their derivations are in \texttt{paper/data/README.md}.

The review log behind Figure~\ref{fig:rounds} is published as release \texttt{review-log-v1} of this repository: \texttt{review-log.tar.gz} (checksum in \texttt{review-log.sha256}) holds the \reviewRoundsTotal{} fix briefs, the re-review protocol that Appendix~\ref{app:protocol} condenses, and the same rounds table, so the fix-round counts repeat from a public artifact~\cite{benchmarkfamily2026reviewlog}. Inside the archive, \texttt{ls fix-briefs | sed -E 's/-fix[0-9]*\textbackslash.md//' | sort | uniq -c} reproduces the per-package round counts. The protocol the first-round defect-finding sessions followed is not part of that release.

\paragraph{Data availability.}
Every number in this paper derives from the data files under \texttt{paper/data/} and the scripts in this repository, and the archive above repeats the review record; the paper itself is archived under concept DOI \texttt{10.5281/zenodo.23428801}. The author's ORCID is \texttt{0009-0005-9964-6841}.

The software is MIT, in this repository and in each benchmark repository. The paper text is released under CC BY 4.0. npm packages and GitHub repositories are the artifact. Each package has a Zenodo concept DOI, listed in its \texttt{CITATION.cff} and README.

\paragraph{Licences per source.}
The tasks are synthetic or written by the maintainer, except the pelican-on-a-bicycle prompt popularised by Simon Willison and the OpenStax formula citation in skateboard; the pelican package declares CC0 for the canonical prompt and MIT for the seeded grid it generates, and the prompt's licence was not verified independently. Table~\ref{tab:suite} lists the licences per package: code under MIT, the seeded grids and catalogs under MIT or CC0, the skateboard glossary citing CC-BY material, and the OpenRouter price feed behind \texttt{cost-per-task} carrying \texttt{NOASSERTION} because the endpoint states no SPDX licence. The paper text is CC BY 4.0.

The directory at \url{https://private-repo-337e.com} lists the packages; it publishes no model ranking.
